\chapter{Overview}

\textbf{HugMesh Games} is a browser-native multi-game portal: a single
Angular shell loads multiple Godot 4 WebAssembly titles on demand, with a
Spring Boot BFF streaming game state and controls over Server-Sent Events.
Authentication runs through Casdoor (OAuth2/OIDC) behind a federated login
remote.

\section{Three sentences}

The host shell is a \textbf{native-federation} Angular host that loads its
\texttt{login}, \texttt{admin}, and per-game routes from sibling bundles
served by per-remote nginx containers on the compose network. The BFF
compiles to either a fat JAR or a \textbf{GraalVM native image}, exposes
the SSE stream, the federation manifest, the upload pipeline, and the
per-game launch / access endpoints, and persists state in PostgreSQL with
a Cassandra cross-pod cache. Each Godot game lives in its own project
under \snapfile{games/<slug>/}, exports a self-contained WASM bundle, and
is delivered to the browser either through the BFF asset path or through
the dedicated \texttt{wasm-fragment} nginx remote with the
COOP/COEP/CORP headers WASM games need.

\section{What's in the box}

\begin{itemize}
  \item \textbf{Angular host} (\reposrc{frontend/}) — federation host plus
        four sibling remotes (\texttt{login}, \texttt{admin},
        \texttt{games-three-wizards}, \texttt{games-dodecahedron}) and
        the asset-tier \texttt{wasm-fragment}.
  \item \textbf{Spring Boot BFF} (\reposrc{backend/}) — WebFlux + R2DBC,
        Liquibase via JDBC, Casdoor JWT verification, GraalVM-ready.
  \item \textbf{Per-game Godot projects} (\reposrc{games/}) — currently
        \texttt{three-wizards-in-a-trenchcoat} (cave + field) and
        \texttt{dodecahedron-demo} (chrome-less single-scene tech demo).
  \item \textbf{Reusable Rust GDExtension} (\reposrc{extensions/hugmesh-bridge/})
        — native \texttt{.so} for the editor, \texttt{wasm32-unknown-emscripten}
        for the web build, with a typed signal bus that any new game can
        consume.
  \item \textbf{wasm-fragment nginx} (\reposrc{wasm-fragment/}) — bakes
        \snapfile{backend/src/main/resources/assets/levels/} and serves
        Godot bundles at \texttt{/\_remotes/wasm-fragment/<slug>/index.html}
        with COOP \texttt{same-origin}, COEP \texttt{require-corp}, and
        CORP \texttt{cross-origin}.
  \item \textbf{Casdoor} (\reposrc{casdoor/}) — never exposed to the host;
        the BFF reverse-proxies the login flow so the admin UI stays
        internal.
  \item \textbf{PostgreSQL + Cassandra} — Postgres holds the canonical
        relational state; Cassandra is the cross-pod movement / session
        cache and a write-through audit log.
\end{itemize}

\section{The plan that produced this shape}

The post-federation topology landed under
\reposrc{docs/PLAN-module-federation-and-native.md}, an eight-section plan
spanning module federation (A), Godot editor plugin (B), hot upload (C),
GraalVM native-image (D), JWT roles + per-game access (E), DB-backed
federation discovery (F), per-game project layout (G), and cryptographic
launch UIDs (H). Every section is at a deliberate end state at the head
of the \reposrc{feature/module-federation-and-native} branch; this
document is the printed companion that explains why each piece exists and
how they fit.

\begin{snapnote}[title={Audience}]
This reference is for engineers who need to extend, deploy, or debug the
portal. The administrator manual, the entry-level tour, and the Godot
developer walkthrough are separate slide decks under
\reposrc{docs/publishing/slides/} and target their respective audiences
without overlap.
\end{snapnote}

\section{What's new since the post-federation snapshot}

Several feature batches have landed on top of the eight-section
post-federation plan; each is covered in detail in its own chapter:

\begin{itemize}
  \item \textbf{Multiplayer sessions.} Per-session SSE sinks
        (\snapfile{ConcurrentHashMap<String, SessionSinks>}) layered over
        the legacy global sink; \texttt{applyControl(event, characterId,
        sessionId)} is the multiplayer-aware overload.
        \snapfile{SessionIsolationTest} pins the cross-session
        non-leakage contract. See chapter~\ref{ch:runtime}.
  \item \textbf{Friends, blocks, session visibility.} Single
        \texttt{user\_relationships} graph table backs both mutual
        friendship and one-directional blocks; \texttt{game\_sessions
        .visibility} is one of \texttt{public} / \texttt{friends} /
        \texttt{private}. The lobby filter runs against the caller's
        friend graph + blocklist. Decision 015 has the full design.
  \item \textbf{Mod and publisher roles with per-game scope.} Casdoor
        owns role membership; Postgres owns the per-game scope in
        \texttt{mod\_game\_scopes} and \texttt{publisher\_game\_scopes}.
        See chapter~\ref{ch:datamodel}.
  \item \textbf{RPG + VN makers.} Two reactive Spring Boot side services
        (\reposrc{rpg-builder/} and \reposrc{vn-builder/}) with their
        own Postgres instances, federated frontend remotes at
        \texttt{/rpg-maker} and \texttt{/vn-maker}, and a build pipeline
        that POSTs the produced bundle to the BFF's
        \texttt{/api/admin/games/upload} endpoint with the publisher's
        bearer. Decision 018 covers the side-service architecture.
  \item \textbf{Production rate-limit zones.} Six nginx
        \texttt{limit\_req\_zone} declarations in
        \snapfile{frontend/nginx.conf} (and the TLS twin) sized for
        $\sim$10 000 concurrent users. The BFF still runs its per-pod
        \snapfile{RateLimitFilter} + Cassandra cross-pod aggregator
        behind the edge. Decision 019.
  \item \textbf{LAN access + TLS overlay.} \snapfile{lan.properties}
        env-file with literal \texttt{FRONTEND\_BIND=0.0.0.0:} values;
        \texttt{make lan-up\{,-tls\}} wraps the stack with self-signed
        TLS (cert SAN includes detected LAN IPs).
  \item \textbf{Crypto.mintUid() + Liquibase 036.} 43-char URL-safe
        base64 SHA3-256 minter replacing \texttt{UUID.randomUUID()} for
        every UI-facing UID; Liquibase 036 widens session-id columns to
        \texttt{varchar(64)} to fit the new shape without disturbing
        pre-existing UUID rows.
\end{itemize}
